Đúng, tôi hiểu ý bạn. Và cách bạn phát biểu này **chính xác hơn cách tôi vừa mô tả trước đó**, vì phải phân biệt rõ ba tầng:

* **Class metamodel**: quy định thế nào là một class diagram hợp lệ.
* **Model \(M\)**: một class diagram cụ thể, gồm classes, attributes, relations, inheritance **và tập OCL \(O\)**.
* **Snapshot \(S\)**: một trạng thái instance của \(M\), tức các object/link cụ thể và phải conform to \(M\).

Khi đó bài toán của bạn là thế này.

$$
MM_{Class}
$$

$$
M \models MM_{Class}
$$

và:

$$
S \models M.
$$

Trong \(M\) có:

$$
O(M)=\{o_1,\ldots,o_n\}.
$$

Từ \(M\), tool sinh tập query:

$$
\boxed{
Compile(M)
=
Q
=
\{q_1,\ldots,q_n\}
}
$$

trong đó:

$$
q_i = Compile(M,o_i).
$$

Còn từ snapshot \(S\), tool sinh property graph:

$$
\boxed{
Build(M,S)=G.
}
$$

Điểm quan trọng là **query được sinh từ model/OCL**, còn **graph được sinh từ snapshot**.

Đó mới là hai luật chuyển của bạn.

---

## Hai nhánh của hệ thống

Ta có:

```text
                  Model M
          (Classes + Relations + OCL)
             /                 \
            /                   \
       OCL compilation        Snapshot S
            |                  S conforms M
            |                       |
            v                       v
       Queries Q                Build
                                    |
                                    v
                               Graph G
```

Hay chính xác hơn:

```text
M = (Class Diagram, OCL set O)
│
├────────────── Compile ──────────────► Q = {q1,...,qn}
│
│
└─ Snapshot S, S conforms M
                 │
                 │ Build
                 ▼
                 G
```

Sau đó correctness mới được xét.

---

# Đối với mỗi OCL \(o_i\)

Giả sử invariant:

$$
o_i\in O(M)
$$

có context class \(K_i\).

Khi đánh giá OCL trực tiếp trên snapshot:

$$
Eval_{OCL}(o_i,S)
$$

ta không chỉ lấy một Boolean.

Ta lấy **tập các object của snapshot vi phạm invariant**:

$$
\boxed{
Viol_S(o_i)
=
\{
s\in Obj(S)
\mid
type(s)\leq K_i
\land
Eval_{OCL}(o_i,s,S)\neq true
\}.
}
$$

Ví dụ:

```ocl
context Person inv PositiveAge:
    self.age > 0
```

Snapshot có:

```text
p1.age = 20
p2.age = -3
p3.age = 50
```

thì:

$$
Viol_S(o)=\{p_2\}.
$$

---

## Bên graph

Snapshot:

$$
S
$$

được chuyển thành:

$$
G=Build(M,S).
$$

Với mỗi object \(s\in S\), phải có một node tương ứng trong \(G\).

Ta cần một ánh xạ:

$$
\boxed{
\gamma_S:Obj(S)\rightarrow ObjNodes(G)
}
$$

và quan trọng:

$$
\boxed{
\gamma_S
\text{ là song ánh}.
}
$$

Tức là:

### mỗi object có đúng một node

$$
\forall s\in Obj(S).
\exists!v\in ObjNodes(G):
v=\gamma_S(s)
$$

và mỗi object-node trong graph đến từ đúng một source object.

Không mất object:

$$
s_1\neq s_2
\Rightarrow
\gamma_S(s_1)\neq\gamma_S(s_2).
$$

Không sinh object-node ma:

$$
\forall v\in ObjNodes(G).
\exists s\in Obj(S):
v=\gamma_S(s).
$$

Đây chính xác là điều bạn đang nói “mỗi \(s\) phải ánh xạ 1-1 với bên graph”.

---

# Query tương ứng

Invariant:

$$
o_i
$$

được compiler sinh:

$$
q_i=Compile(M,o_i).
$$

Chạy:

$$
q_i(G)
$$

trả về tập node vi phạm:

$$
\boxed{
Ans_G(q_i)\subseteq ObjNodes(G).
}
$$

Correctness mà bạn cần chứng minh là:

$$
\boxed{
Ans_G(q_i)
=
\gamma_S[Viol_S(o_i)].
}
\tag{1}
$$

Đây mới là câu quan trọng nhất.

Nghĩa là nếu:

$$
Viol_S(o_i)=\{s_2,s_5,s_8\}
$$

thì query phải trả đúng:

$$
Ans_G(q_i)
=
\{
\gamma(s_2),
\gamma(s_5),
\gamma(s_8)
\}.
$$

Không hơn.

Không thiếu.

---

# Và theorem của toàn phương pháp

Tôi sẽ phát biểu theorem của bạn như sau:

$$
\boxed{
\begin{aligned}
\forall M,S.\quad
&
Conforms_{MM}(M)
\land
Conforms_M(S)
\\
&\Longrightarrow
\\
&
\exists\gamma_S.
\ Bijective(\gamma_S)
\\
&\qquad{}\land
\forall o\in O(M).
\\
&
\gamma_S[Viol_{OCL}(S,o)]
=
Ans_{Cypher}
\left(
Build(M,S),
Compile(M,o)
\right).
\end{aligned}
}
\tag{MASTER}
$$

Đây là cách tôi nghĩ **đúng nhất với case study bạn vừa mô tả**.

---

# Có thể viết ngắn hơn

Đặt:

$$
G = Build(M,S)
$$

và:

$$
Q=Compile(M,O(M)).
$$

Có ánh xạ:

$$
\gamma:Obj(S)\cong ObjNodes(G).
$$

Thì correctness là:

$$
\boxed{
\forall o_i\in O(M):
\quad
\gamma
\left[
EvalViol_{OCL}(o_i,S)
\right]
=
EvalViol_{Graph}
(q_i,G).
}
\tag{2}
$$

với:

$$
q_i=Compile(M,o_i).
$$

Đây chính là **semantic preservation** của hai transformation phối hợp.

---

# Bây giờ hình vuông correctness mới chuẩn

Bạn có:

```text
                         evaluate OCL oi
       Snapshot S  ─────────────────────────►  ViolS(oi)
           │                                      │
           │ Build(M,S)                           │ γ
           │                                      │
           ▼                                      ▼
       Graph G      ─────────────────────────► AnsG(qi)
                        execute qi
```

với:

$$
q_i=Compile(M,o_i).
$$

Bạn phải chứng minh hình vuông commute:

$$
\boxed{
\gamma\circ Eval_{OCL}
=
Eval_{Cypher}\circ(Build,Compile).
}
$$

Hay trên tập violation:

$$
\boxed{
\gamma[Viol_{OCL}(S,o)]
=
Viol_{Cypher}(Build(M,S),Compile(M,o)).
}
$$

Đây là formulation rất đẹp cho paper.

---

# Vậy Isabelle phải biểu diễn theo formulation mới này

Điểm này cũng làm tôi sửa lại kiến trúc Isabelle trước đó.

Isabelle cần có trước tiên:

```text
ClassMetamodel
      │
      │ conforms
      ▼
    Model M
(Class/Attr/Relation/OCL)
      │
      │
      ├──────────── Compile ───────► Query
      │
      ▼
  Snapshot S
      │ conforms M
      │
      │ Build
      ▼
  PropertyGraph G
```

Nghĩa là trong Isabelle phải phân biệt rõ:

```isabelle
class_metamodel
```

```isabelle
uml_model
```

```isabelle
snapshot
```

chứ **không được gộp `uml_model` và `snapshot` như cách tôi nói trước**.

Sau đó:

```isabelle
conforms_model ::
  "class_metamodel ⇒ uml_model ⇒ bool"
```

và:

```isabelle
conforms_snapshot ::
  "uml_model ⇒ snapshot ⇒ bool"
```

---

## Model \(M\)

Có thể conceptual:

```isabelle
record uml_model =
  classes      :: ...
  attributes   :: ...
  associations :: ...
  inheritance  :: ...
  invariants   :: "ocl_inv set"
```

---

## Snapshot \(S\)

```isabelle
record snapshot =
  objects :: "object set"
  links   :: "link set"
```

Snapshot không chứa class definitions.

Nó là instance:

$$
S\models M.
$$

---

# Transformation 1

```isabelle
compile ::
  "uml_model ⇒ ocl_inv ⇒ query"
```

hoặc toàn model:

```isabelle
compile_all ::
  "uml_model ⇒ (ocl_inv × query) set"
```

với:

$$
Compile(M)
=
\{(o,Compile(M,o))\mid o\in O(M)\}.
$$

---

# Transformation 2

```isabelle
build ::
  "uml_model ⇒ snapshot ⇒ property_graph"
```

Tôi nghĩ signature này chính xác hơn:

$$
\boxed{
Build:M\times S\rightarrow G.
}
$$

Vì để chuyển snapshot sang graph, builder phải biết structure của \(M\):

* class nào;
* association nào;
* multiplicity;
* inheritance;
* role;
* encoding.

Không nên chỉ viết:

$$
Build(S).
$$

Mà nên là:

$$
\boxed{
Build(M,S).
}
$$

---

# Sau đó định nghĩa violation phía nguồn

```isabelle
ocl_violations ::
  "uml_model ⇒ snapshot ⇒ ocl_inv ⇒ object set"
```

$$
Viol_O(M,S,o).
$$

Và phía graph:

```isabelle
query_violations ::
  "property_graph ⇒ query ⇒ node set"
```

$$
Viol_Q(G,q).
$$

---

# Cuối cùng là correspondence

```isabelle
rep ::
  "snapshot ⇒ property_graph ⇒ object ⇒ node"
```

hoặc existential witness:

$$
\gamma:Obj(S)\leftrightarrow Nodes(G).
$$

Rồi theorem:

```isabelle
theorem transformation_correct:
  assumes "conforms_model ClassMM M"
      and "conforms_snapshot M S"
      and "o ∈ invariants M"
  defines "G ≡ build M S"
  defines "q ≡ compile M o"
  shows
    "∃γ.
       bij_betw γ (objects S) (object_nodes G)
       ∧ γ ` ocl_violations M S o
         = query_violations G q"
```

**Đây gần như chính xác lời bạn vừa nói, nhưng viết thành Isabelle theorem.**

---

Và tôi nghĩ có một insight rất quan trọng ở đây: **cái cần prove không phải riêng “OCL→Cypher đúng” hay riêng “Snapshot→Graph đúng”. Cái cần prove là hai phép chuyển phối hợp sao cho tập object vi phạm được bảo toàn qua ánh xạ object→node.**

Đó mới thực sự là correctness property của case study của bạn.
